\paragraph{The model (Definition 1).}
We say that $f:W^N\to\Sigma^m$ is a depth-$q$ forest with word size $w$ and output alphabet $\Sigma$ if $W=\{0,1\}^w$ and $f=(f_1,f_2,\ldots,f_m)$ where each $f_i:W^N\to\Sigma$ is a depth-$q$ decision tree where the variables are over $W$, each internal node has $|W|$ children, and the leaves are labeled with elements from $\Sigma$.
The main goal of this paper is to show that a target distribution $S$ over $\Sigma^m$ is hard to sample by a low-depth forest. We measure the distance between distributions $X$ and $Y$ over $D$ using statistical (a.k.a. total variation, $L_1$) distance
\[\Delta(X,Y):=\max_{T\subseteq D}|\mathbb P[X\in T]-\mathbb P[Y\in T]|.\]
So the lower-bound goal is to show $\Delta(f(U_{W^N}),S)$ is large for any low-depth forest $f$, where $U_{W^N}$ is the uniform distribution over $W^N$. In general for a set $H$ we write $U_H$ for the uniform distribution over $H$, and simply $U$ when $H$ is clear from the context.
\paragraph{Definition 13 (Almost pairwise independence).}
Jointly distributed random variables $X,Y$ are $\epsilon$-independent if $(X,Y)$ is $\epsilon$-close in statistical distance to $(X,Y')$ where $Y'$ has the same distribution of $Y$ and is independent from $X$.
The min-entropy $H_\infty(X)$ of a random variable $X$ is $\min_a\log_2(1/\mathbb P[X=a])$.
\paragraph{Notation.}
To avoid clutter in the more technical exposition of the results, we adopt the convention that for a set $S$ we also denote by $S$ the uniform distribution $U_S$ over $S$. The meaning will be clear from the context. For example, we shall simply write $\Delta(f(W^N),S)$ for $\Delta(f(U_{W^N}),U_S)$.
